\subsection{Alignment of Twisted Spectral Gaps: Unification Criteria and Group Selectors}\label{subsec:group-unification-spectral-alignment}

\paragraph{Spectral gap as the endogenous location of channel scales}
In the kernel Chebotarev estimate, the leading term of the congruence class stratification of primitive orbits is controlled by the Perron root \(\lambda_1(u)\),
while the exponential rate of the tail term is controlled by the worst-case twisted spectral radius (and the second-largest modulus eigenvalue of the untwisted case) (Theorem~\ref{thm:kernel-chebotarev-exp}).
Therefore, the spectral gap ratio \(\Lambda(u)/\lambda_1(u)\) and its logarithm
\(
g(u)=\log(\lambda_1(u)/\Lambda(u))
\)
constitute the auditable location of the ``readout time scale/mixing time scale'' (Corollary~\ref{cor:pom-order-mixing} and the same-type caliber of Proposition~\ref{prop:pom-rewrite-audit-loop}).

\begin{definition}[Spectral coupling]\label{def:spec-coupling}
In the notation of Definition~\ref{def:kernel-group-twist}, fix the temperature parameter \(u\in(0,1]\) and write \(\lambda_1(u)=\rho(M_K(u))\).
For any irreducible complex representation \(\varrho\in\widehat G\), write its twisted operator spectral radius
\[
\eta_{\varrho}(u):=\rho\!\bigl(M_{K,\varrho}(u)\bigr).
\]
Define the spectral gap coupling of the representation channel:
\[
g_{\varrho}(u):=\log\!\left(\frac{\lambda_1(u)}{\eta_{\varrho}(u)}\right).
\]
Also write the worst-case nontrivial twisted scale:
\[
g_{\mathrm{tw}}(u):=\log\!\left(\frac{\lambda_1(u)}{\eta(u)}\right)
=\min_{\varrho\in\widehat G,\ \varrho\ne \mathbf{1}} g_{\varrho}(u),
\qquad
\eta(u):=\max_{\varrho\ne\mathbf{1}}\eta_{\varrho}(u)
\]
(consistent with the \(\eta(u)\) of Theorem~\ref{thm:kernel-chebotarev-exp}).
\end{definition}

\begin{remark}[Minimax unification criterion for spectral gap alignment]\label{rem:spec-gap-minimax-unification}
Let \(\mathcal{R}_{\mathrm{obs}}\subseteq \widehat G\) be a family of ``observable channel'' representations (e.g., the twisted coordinate families activated by specific readout gates).
Write the discrete mean of \(\{g_{\varrho}(u)\}_{\varrho\in\mathcal{R}_{\mathrm{obs}}}\) as
\(
\bar g(u):=|\mathcal{R}_{\mathrm{obs}}|^{-1}\sum_{\varrho\in\mathcal{R}_{\mathrm{obs}}}g_{\varrho}(u)
\).
Define the worst-case deviation function for spectral gap alignment:
\[
\mathcal{A}(u):=\max_{\varrho\in\mathcal{R}_{\mathrm{obs}}}\left|g_{\varrho}(u)-\bar g(u)\right|.
\]
Then ``unification'' can be stated as a purely spectral-theoretic minimax problem that is enumerable/auditable at the finite matrix level:
over the allowed kernel parameters and weight slices, find the unification point \(u_\ast\) that minimizes \(\mathcal{A}(u)\).
\par
This caliber does not require prior commitment to any continuous gauge group; the input objects for unification are the computable spectral data of the twisted matrix family \(M_{K,\varrho}(u)\),
and the output object is the testable degree of spectral gap alignment (quantified by \(\mathcal{A}(u)\)).
\end{remark}

\paragraph{Structural Consequences of the Unification Functional \texorpdfstring{$\mathcal{A}(u)$}{A(u)}: Exact Cancellation, Local Optimality, and Algebraic Certificates}\label{par:gut-unification-functional-structures}
This paragraph, without repeating the premises defined in Remark \ref{rem:spec-gap-minimax-unification},
extracts within its existing notation framework several conclusions that have nontrivial consequences for the unification functional \(\mathcal{A}(u)\) itself;
these conclusions are independent of the ``trivial sector'' of \(\lambda_1(u)\) and are suitable as independent mathematical citation units.

\paragraph{Geometric Mean Coordinates and Exact Cancellation of the Trivial Sector}
Following the notation of Remark \ref{rem:spec-gap-minimax-unification}, let
\(
\mathcal{R}:=\mathcal{R}_{\mathrm{obs}}
\)
and
\(
R:=|\mathcal{R}|
\).
Define the geometric mean and extrema of the twisted spectral radii
\[
\eta_{\mathrm{geo}}(u):=\Bigl(\prod_{\varrho\in\mathcal{R}}\eta_\varrho(u)\Bigr)^{1/R},
\qquad
\eta_{\min}(u):=\min_{\varrho\in\mathcal{R}}\eta_\varrho(u),
\qquad
\eta_{\max}(u):=\max_{\varrho\in\mathcal{R}}\eta_\varrho(u).
\]

\begin{proposition}[Exact Cancellation and Intrinsic Form]\label{prop:gut-A-cancels-lambda1}
For each \(u\in(0,1]\), the following identity holds:
\[
\mathcal{A}(u)
=\max_{\varrho\in\mathcal{R}}\Bigl|\log\eta_{\mathrm{geo}}(u)-\log\eta_\varrho(u)\Bigr|
=\log\max\Bigl\{\frac{\eta_{\mathrm{geo}}(u)}{\eta_{\min}(u)},\ \frac{\eta_{\max}(u)}{\eta_{\mathrm{geo}}(u)}\Bigr\}.
\]
In particular, \(\mathcal{A}(u)\) and its minimizer are completely \emph{insensitive} to \(\lambda_1(u)\): although each \(g_\varrho(u)\) contains \(\lambda_1(u)\), its contribution cancels exactly after centering.
\end{proposition}

\begin{proof}
By definition \(g_\varrho(u)=\log\lambda_1(u)-\log\eta_\varrho(u)\), hence
\[
\bar g(u)=\log\lambda_1(u)-\frac{1}{R}\sum_{\sigma\in\mathcal{R}}\log\eta_\sigma(u),
\]
so that
\[
g_\varrho(u)-\bar g(u)
=\frac{1}{R}\sum_{\sigma\in\mathcal{R}}\log\eta_\sigma(u)-\log\eta_\varrho(u)
=\log\eta_{\mathrm{geo}}(u)-\log\eta_\varrho(u).
\]
Taking the absolute value and maximizing over \(\varrho\) yields the first expression.
The second expression follows immediately from \(\eta_{\min}\le \eta_{\mathrm{geo}}\le \eta_{\max}\). \qed
\end{proof}

\begin{corollary}[Two-Channel Collapse]\label{cor:gut-A-two-channel}
If \(R=2\), then
\[
\mathcal{A}(u)=\frac12\log\frac{\eta_{\max}(u)}{\eta_{\min}(u)}.
\]
Therefore in the two-channel case, minimizing \(\mathcal{A}(u)\) is completely equivalent to minimizing the logarithmic spread \(\log(\eta_{\max}/\eta_{\min})\).
\end{corollary}

\begin{corollary}[Spread Equivalence (up to a Constant Factor)]\label{cor:gut-A-spread-equivalence}
For all \(R\ge 2\) and \(u\in(0,1]\), the following two-sided bound holds:
\[
\frac12\log\frac{\eta_{\max}(u)}{\eta_{\min}(u)}
\le
\mathcal{A}(u)
\le
\log\frac{\eta_{\max}(u)}{\eta_{\min}(u)}.
\]
Therefore \(\mathcal{A}(u)\) is equivalent to the logarithmic condition number of the multiset \(\{\eta_\varrho(u)\}_{\varrho\in\mathcal{R}}\) (up to a factor of \(2\)).
\end{corollary}

\paragraph{Intrinsic \texorpdfstring{$\beta$}{beta}-Field and Necessary and Sufficient Criteria for Local Optimality}
On the scale \(\theta=\log u\), let (where the derivative exists)
\[
b_\varrho(u):=\frac{\dd}{\dd\log u}\log\eta_\varrho(u),
\qquad
b_{\mathrm{geo}}(u):=\frac{1}{R}\sum_{\sigma\in\mathcal{R}}b_\sigma(u)
=\frac{\dd}{\dd\log u}\log\eta_{\mathrm{geo}}(u).
\]
Define also
\[
h_\varrho(u):=\log\eta_{\mathrm{geo}}(u)-\log\eta_\varrho(u),
\qquad
\mathcal{A}(u)=\max_{\varrho\in\mathcal{R}}|h_\varrho(u)|
\]
(by Proposition \ref{prop:gut-A-cancels-lambda1}).
Furthermore, let \(b_1(u):=\frac{\dd}{\dd\log u}\log\lambda_1(u)\), then
\(
\beta_\varrho(u)=b_1(u)-b_\varrho(u)
\)
and after centering
\(
\beta_\varrho(u)-\bar\beta(u)=-(b_\varrho(u)-b_{\mathrm{geo}}(u))
\);
hence the first-order variation of \(\mathcal{A}(u)\) depends only on the intrinsic centered quantities \(b_\varrho-b_{\mathrm{geo}}\).

\begin{theorem}[Necessary and Sufficient Conditions for Local Minimum (One-Dimensional KKT/Clarke Conditions for Finite Maxima)]\label{thm:gut-A-local-optimality}
Suppose each \(u\mapsto \log\eta_\varrho(u)\) is differentiable at \(u_\star\in(0,1]\).
Define the active set and signs
\[
\mathcal{A}_{\mathrm{act}}(u_\star):=\Bigl\{\varrho\in\mathcal{R}:\ |h_\varrho(u_\star)|=\mathcal{A}(u_\star)\Bigr\},
\qquad
\sigma_\varrho:=\operatorname{sgn}\bigl(h_\varrho(u_\star)\bigr)\in\{\pm 1\}.
\]
Then \(u_\star\) is a local minimum of \(\mathcal{A}\) if and only if
\[
0\in
\operatorname{conv}\Bigl\{\sigma_\varrho\bigl(b_{\mathrm{geo}}(u_\star)-b_\varrho(u_\star)\bigr):\ \varrho\in\mathcal{A}_{\mathrm{act}}(u_\star)\Bigr\}.
\]
Equivalently, there exist coefficients \(\alpha_\varrho\ge 0\) (\(\varrho\in\mathcal{A}_{\mathrm{act}}(u_\star)\)) satisfying
\(
\sum_{\varrho\in\mathcal{A}_{\mathrm{act}}(u_\star)}\alpha_\varrho=1
\)
and
\[
\sum_{\varrho\in\mathcal{A}_{\mathrm{act}}(u_\star)}
\alpha_\varrho\,\sigma_\varrho\,\bigl(b_{\mathrm{geo}}(u_\star)-b_\varrho(u_\star)\bigr)=0.
\]
In the typical ``dual active channel'' case (\(\mathcal{A}_{\mathrm{act}}(u_\star)=\{\varrho_+,\varrho_-\}\) with \(h_{\varrho_+}(u_\star)=\mathcal{A}(u_\star)\) and \(h_{\varrho_-}(u_\star)=-\mathcal{A}(u_\star)\)), the above condition is equivalent to
\[
\bigl(b_{\mathrm{geo}}(u_\star)-b_{\varrho_+}(u_\star)\bigr)\cdot
\bigl(b_{\mathrm{geo}}(u_\star)-b_{\varrho_-}(u_\star)\bigr)\le 0.
\]
\end{theorem}

\begin{proof}
Let \(\theta=\log u\) and write \(F(\theta):=\mathcal{A}(e^\theta)\).
By the definition of \(h_\varrho\), \(F\) can be expressed in a neighborhood of \(\theta_\star=\log u_\star\) as the maximum of finitely many differentiable functions:
\(
F(\theta)=\max_{\varrho\in\mathcal{R}}\max\{h_\varrho(e^\theta),-h_\varrho(e^\theta)\}.
\)
In the one-dimensional case, a finite maximum function has a local minimum at \(\theta_\star\) if and only if \(0\) belongs to the convex hull of the derivatives of its active branches (the KKT criterion for the Clarke subdifferential).
Computing the derivative of the active branch \(\sigma_\varrho h_\varrho(e^\theta)\) gives
\[
\frac{\dd}{\dd\theta}\bigl(\sigma_\varrho h_\varrho(e^\theta)\bigr)\Big|_{\theta=\theta_\star}
=\sigma_\varrho\Bigl(\frac{\dd}{\dd\log u}\log\eta_{\mathrm{geo}}(u)-\frac{\dd}{\dd\log u}\log\eta_\varrho(u)\Bigr)\Big|_{u=u_\star}
=\sigma_\varrho\bigl(b_{\mathrm{geo}}(u_\star)-b_\varrho(u_\star)\bigr),
\]
yielding the stated necessary and sufficient condition. In the dual active channel case, the convex hull condition is equivalent to the two numbers having opposite signs or containing zero, i.e., their product being non-positive. \qed
\end{proof}

\paragraph{Stability of the Unification Point under Uniform Error: From \texorpdfstring{$\|\widetilde{\mathcal A}-\mathcal A\|_\infty$}{sup} to Argmin Error Bounds}
\begin{theorem}[Minimizer Stability under Uniform Error]\label{thm:gut-argmin-stability}
Let \(I\subset\RR\) be a compact interval (with variable \(\theta=\log u\)), and write \(F(\theta):=\mathcal{A}(e^\theta)\).
Assume \(F\) has a unique minimizer \(\theta_\star\) on \(I\), satisfying the sharpness inequality
\[
F(\theta)\ge F(\theta_\star)+\kappa|\theta-\theta_\star|\qquad(\forall\,\theta\in I)
\]
where \(\kappa>0\).
If an approximate function \(\widetilde F\) satisfies \(\|\widetilde F-F\|_{L^\infty(I)}\le \varepsilon\), then any minimizer \(\widetilde\theta_\star\) of \(\widetilde F\) on \(I\) satisfies
\[
|\widetilde\theta_\star-\theta_\star|\le \frac{2\varepsilon}{\kappa},
\qquad
F(\widetilde\theta_\star)\le F(\theta_\star)+2\varepsilon.
\]
\end{theorem}

\begin{proof}
By the optimality of \(\widetilde\theta_\star\) and the uniform error bound,
\[
F(\widetilde\theta_\star)\le \widetilde F(\widetilde\theta_\star)+\varepsilon
\le \widetilde F(\theta_\star)+\varepsilon
\le F(\theta_\star)+2\varepsilon.
\]
Applying the sharpness inequality at \(\theta=\widetilde\theta_\star\) yields
\(
\kappa|\widetilde\theta_\star-\theta_\star|\le F(\widetilde\theta_\star)-F(\theta_\star)\le 2\varepsilon
\),
and the conclusion follows. \qed
\end{proof}

\paragraph{Determinant Elimination Certificate: Replacing Numerical Differentiation with \texorpdfstring{$\det(I-zM(u))$}{det}}
In the core audit caliber of this paper, many channel spectral radii arise from the Perron root of finite-dimensional matrix families \(M_{K,\varrho}(u)\);
therefore the intrinsic \(\beta\)-field \(b_\varrho(u)=\frac{\dd}{\dd\log u}\log\eta_\varrho(u)\) can be fully algebraized via partial derivative ratios of the determinant.

\begin{proposition}[Logarithmic Derivative of the Spectral Radius via Partial Derivative Ratio]\label{prop:gut-beta-determinant-ratio}
Let \(M(u)\) be a finite matrix whose entries are rational functions of \(u\), and suppose on some connected open interval \(U\subset(0,1]\) there exists a differentiable function \(z_\star(u)\in\RR_{>0}\) such that
\[
\Delta(z,u):=\det(I-zM(u))=0\quad\text{at}\quad z=z_\star(u)
\]
with \(z_\star(u)^{-1}=\rho(M(u))\), and this root is simple on \(U\) (i.e., \(\partial_z\Delta(z_\star(u),u)\neq 0\)).
Let \(\eta(u):=\rho(M(u))=z_\star(u)^{-1}\) and define
\(
b(u):=\frac{\dd}{\dd\log u}\log\eta(u)
\).
Then for all \(u\in U\) the following identity holds:
\[
b(u)=\frac{u}{z_\star(u)}\,
\frac{\partial_u\Delta(z_\star(u),u)}{\partial_z\Delta(z_\star(u),u)}.
\]
Moreover, after clearing denominators of \(\Delta\) (so that it lies in some \(\ZZ[u^{\pm1},z]\)), \((b,u)\) satisfies an elimination relation in \(z\): there exists a nonzero polynomial
\(
\mathcal{R}(b,u)\in\ZZ[u^{\pm1},b]
\)
such that \(\mathcal{R}(b(u),u)=0\) (which can be taken as the resultant with respect to \(z\)).
\end{proposition}

\begin{proof}
Differentiating the identity \(\Delta(z_\star(u),u)=0\) with respect to \(u\) yields
\[
\partial_z\Delta(z_\star(u),u)\,z_\star'(u)+\partial_u\Delta(z_\star(u),u)=0,
\]
and by the simple root condition,
\(
z_\star'(u)=-\partial_u\Delta/\partial_z\Delta
\)
(evaluated at \(z=z_\star(u)\)).
Since
\(
b(u)=\frac{\dd}{\dd\log u}\log(z_\star(u)^{-1})=-\frac{u\,z_\star'(u)}{z_\star(u)}
\),
substitution gives the first formula.
The second part is a direct consequence of elimination: at \(z=z_\star(u)\), both
\(
\Delta(z,u)=0
\)
and
\(
b\,z\,\partial_z\Delta(z,u)-u\,\partial_u\Delta(z,u)=0
\)
hold simultaneously, and taking the resultant with respect to \(z\) yields an algebraic relation involving only \((b,u)\). \qed
\end{proof}



\begin{remark}[Block spectral gap alignment: strengthening the selector with \texorpdfstring{$9+12$}{9+12} and \texorpdfstring{$1+8+12$}{1+8+12}]\label{rem:spec-gap-block-selector}
If the labels of the observable channel family \(\mathcal{R}_{\mathrm{obs}}\) can be aligned with the boundary tower word table (e.g., labeled by the discrete labels of \(X_6\)),
then \textbf{stratified spectral gap alignment} tests can be performed on the \(9+12\) (\(\mathfrak{u}(3)\) compact core vs.\ noncompact complement, \S\ref{subsec:bdry-tower-zeck-gut})
and \(1+8+12\) (\(SU(5)\) light/heavy dictionary) partitions.
Write \(B_9,B_{12}\) for the \(9+12\) blocks and \(B_1,B_8,B_{12}\) for the \(1+8+12\) blocks.
For any block \(B\subseteq \mathcal{R}_{\mathrm{obs}}\), write
\[
\bar g_B(u):=\frac{1}{|B|}\sum_{\varrho\in B} g_{\varrho}(u),\qquad
\mathcal{A}_B(u):=\max_{\varrho\in B}\left|g_{\varrho}(u)-\bar g_B(u)\right|.
\]
Then the falsifiable signatures at the unification point \(u_\ast\) include:
(i)~if \(\mathcal{A}_{12}(u_\ast) < \mathcal{A}_{9}(u_\ast)\) with positive margin, this supports the \(\mathrm{Sp}(6)\) Cartan decomposition \(9+12\) as the candidate partition of observable channels;
(ii)~if the alignment of the \(1+8\) subblock is better than that of the \(12\) subblock, this supports the \(SU(5)\)-type light/heavy dictionary ``\(21-(8+1)=12\).''
Moreover, using the entropy weight \(h(w)=\log\Xi(w)\) defined above,
a weighted average can be introduced for the \(8\oplus 4\) two-tier decomposition of the heavy remainder \(12\):
\[
\bar g_B^{(h)}(u):=\frac{\sum_{\varrho\in B} h(\varrho)\,g_{\varrho}(u)}{\sum_{\varrho\in B} h(\varrho)},
\]
thereby converting ``two-tier heavy spectrum/entropy weight breaking'' into an auditable spectral gap alignment selector.
\end{remark}

\begin{remark}[Symmetric fixedness at the self-dual point \(u=1\) and six-branch locking]\label{rem:selfdual-u1-six-branch-lock}
Under the self-dual completion caliber of the synchronization kernel, the variable change \(u=e^{\theta}\) compresses \(u\leftrightarrow 1/u\) to the central symmetry \(\theta\leftrightarrow -\theta\);
therefore \(u=1\) (i.e., \(\theta=0\)) is the unique self-dual fixed point (see Appendix~\ref{app:kernel-functional-equation}).
At this fixed point, the pressure satisfies \(P(\theta)-\theta/2\) being an even function, so all odd-order cumulants vanish strictly:
\(
P^{(2k+1)}(0)=0
\)
(Proposition~\ref{prop:fluctuation-even}).
\par
More strongly, the elimination certificate of the rate curve carries a sixfold aggregation factor \((u-1)^6\) at the \(\alpha=\tfrac12\) cross-section
(see the generated certificate of Remark~\ref{rem:rate-algebraic-elimination-structure}); its analytic meaning is:
at \(u=1\), the \(6\) nonzero spectral branches of the synchronization kernel simultaneously satisfy tangent locking (Theorem~\ref{thm:sync-spectral-tangent-lock} and Corollary~\ref{cor:sync-spectral-tangent-lock-6}).
Therefore, under the self-dual caliber, pinning the unification point at \(u=1\) is a rigid choice supported jointly by an algebraic certificate and spectral branch structure.
\end{remark}

\begin{remark}[Rank interpretation of the six-branch locking (working hypothesis and falsifiable test)]\label{rem:u1-six-branch-rank3-certificate}
Under the self-dual completion, spectral data naturally pair under the action \(u\leftrightarrow 1/u\);
therefore an empirical organizing principle is: if the unification criterion depends only on an \(r\)-dimensional family of independent channels (equivalently, \(r\) independent \(\theta\) directions/rank),
then the ``number of tangency-locked nonzero spectral branches'' near the central fixed point \(u=1\) should be \(2r\) (given by the self-dual pairing).
Under this caliber, the sixfold aggregation factor \((u-1)^6\) at the \(\alpha=\tfrac12\) cross-section can be viewed as a candidate ``spectral rank certificate'' for \(r=3\):
it is compatible with the \(m=6\) unification seed being locked to rank-\(3\) \(B_3/C_3\) (and the continuous envelope \(\mathrm{Sp}(6)\)),
and more naturally pushes rank-\(4\) \(SU(5)\) to the ``supplementary closure'' level rather than the ``primary spectral unification'' level (counting closure: Corollary~\ref{cor:su5-21-plus-3-closure}).
\par
This principle itself is falsifiable/auditable: while maintaining the self-dual structure, introducing an additional independent twist/potential direction
and repeating the elimination certificate computation, one checks whether the exponent of the \((u-1)\) factor increases from \(6\) to \(8\).
\end{remark}

\begin{remark}[Auditable numerical anchor: spectral gap ordering and low-order R\'enyi near-monospectral dimension signal]\label{rem:spec-gap-numeric-anchors}
Table~\ref{tab:real_input_40_dirichlet_mertens_tensor_anova} gives the worst-case twisted ratio \(\rho/\lambda\) of the real-input 40-state kernel on several three-dimensional residue tensors,
from which the channel spectral gap quantity \(-\log(\rho/\lambda)\) (i.e., the \(1/\tau_{\mathrm{mix}}\) of Remark~\ref{rem:pom-chi-mixing}) is derived.
On the other hand, Table~\ref{tab:fold_collision_renyi_spectrum} gives the low-order R\'enyi fingerprint \(h_q=\log(2^q/r_q)\) of the fold output distribution,
where the closeness of \(h_3\) and \(2h_2\) serves as a testable signal for ``low-order near-monospectral dimension.''
The numerical values of both derived quantities are generated by a single script run:
\input{sections/generated/eq_group_unification_spectral_alignment_diagnostics}
\end{remark}

\begin{remark}[Character field matching criterion: filtering group labels by constant field closure]\label{rem:character-field-matching}
For a fixed kernel \(K\), the various ``finite-part/cumulant'' constants in this paper algebraically fall into a natural number field:
for example, the output potential of the real-input 40-state kernel satisfies \(P^{(k)}(0)\in\mathbb{Q}(\sqrt5)\)
(Proposition~\ref{prop:real-input-40-output-cumulants-closed}; see also the number field stratification summary in Remark~\ref{rem:residue-field-sync-vs-real}).
\par
On the other hand, the matrix entries of the twisted operator \(M_{K,\varrho}(u)\) are generated by the matrix coefficients of the representation \(\varrho\);
therefore the coefficients of its characteristic polynomial belong to the composite of the ``baseline constant field'' and the ``representation coefficient field,'' and its spectral data (including \(\eta_{\varrho}(u)\)) generally introduce new number field inputs.
This suggests an auditable group selection criterion: if one wishes the twisted spectral constants after unification to be as algebraically closed as possible within the existing constant field,
then one may prioritize group label families whose irreducible representations can be realized over that constant field (or whose character value fields contain the constant field without forcing a larger extension).
\par
In the examples of this paper, the slow-mode worst-case twist is closely related to the \(p=5\) axis (see Remark~\ref{rem:arity-335-kappa-poisson} and Table~\ref{tab:real_input_40_arity_335_n2_selection_law_primes}),
and \(\cos(2\pi/5)\in\mathbb{Q}(\sqrt5)\) provides a ``minimal nontrivial twist angle'' anchor cognate with the golden quadratic field;
this makes the closure constraint among ``constant field---twist channel---group label'' amenable to numerically falsifiable interface.
\end{remark}

\begin{remark}[Spectral gap running law on the discrete modulus axis: \texorpdfstring{$p^{-2}$}{p\^{}-2} diffusive scaling of \texorpdfstring{$\Delta(p)$}{Delta(p)}]\label{rem:real-input-40-gap-diffusive-running}
For the real-input 40-state kernel, along the \(((p,p,p))\) residue tensor family (third axis taking \(N_e\bmod p\)), write
\[
\Delta(p):=-\log(\rho_{p,p,p}/\lambda),\qquad
\tau_{\mathrm{mix}}(p):=\frac{1}{\Delta(p)}.
\]
Table~\ref{tab:real-input-40-rho-ppp-prime-scan} gives numerical values at several prime points together with two running diagnostics
\(
C_1(p)=p\,\Delta(p)
\)
and
\(
C_2(p)=p^2\,\Delta(p)
\).
From the table, \(C_2(p)\) stabilizes at a plateau of \(\approx 0.95\text{--}0.96\) for \(p\ge 5\), while \(C_1(p)\) decreases monotonically with \(p\);
therefore this channel spectral gap is more consistent with the diffusive scaling
\[
\Delta(p)\ \asymp\ C\,p^{-2}
\]
rather than a first-order running law \(C\,p^{-1}\).
Correspondingly, the mixing scale grows as \(\tau_{\mathrm{mix}}(p)=\Theta(p^2)\) (consistent with the diffusive mixing caliber of Remark~\ref{cor:arity-335-diffusive-mixing}, but with constants depending on the specific kernel and twist axis choice).
\end{remark}
\input{sections/generated/tab_real_input_40_rho_ppp_prime_scan}

\begin{remark}[Rescaling of mixing time under the external time scaling proxy]\label{rem:lapse-modulated-tau-mix}
If the background-dependent routing/inverse-search overhead is abstracted as a positive quantity \(\kappa\) (units: external cost per internal step), and the external time scaling proxy is defined with reference value \(\kappa_0\) as
\[
N:=\kappa_0/\kappa,
\]
then under the simplified model of ``external time acts only as rescaling,'' there is a deterministic conversion
\[
\Delta_{\mathrm{eff}}=N\Delta,\qquad
\tau_{\mathrm{mix,eff}}=\tau_{\mathrm{mix}}/N.
\]
Table~\ref{tab:real-input-40-lapse-modulated-tau-mix} instantiates this conversion as an auditable prediction table: when \(\kappa\) is instantiated by a specific observation/search cost model, the order-of-magnitude change in the external mixing scale is obtained.
\end{remark}
% AUTO-GENERATED by scripts/exp_real_input_40_lapse_modulated_tau_mix.py
\begin{table}[H]
\centering
\scriptsize
\setlength{\tabcolsep}{6pt}
\renewcommand{\arraystretch}{1.15}
\caption{Overhead-modulated mixing-time rescaling: given $\Delta=-\log(\rho/\lambda)$ and $\tau_{\mathrm{mix}}=1/\Delta$ (internal steps), an external-time rescaling proxy $N=\kappa_0/\kappa$ yields $\Delta_{\mathrm{eff}}=N\Delta$ and $\tau_{\mathrm{mix,eff}}=\tau_{\mathrm{mix}}/N$ (external time).}
\label{tab:real-input-40-lapse-modulated-tau-mix}
\begin{tabular}{r r r r r r}
\toprule
$p$ & $\rho/\lambda$ & $\Delta$ & $\kappa$ & $N$ & $\tau_{\mathrm{mix,eff}}$\\
\midrule
3 & 0.936686340205 & 0.065406801766 & 1.000000000000 & 1.000000000000 & 15.288929790195\\
3 & 0.936686340205 & 0.065406801766 & 2.000000000000 & 0.500000000000 & 30.577859580389\\
3 & 0.936686340205 & 0.065406801766 & 10.000000000000 & 0.100000000000 & 152.889297901946\\
3 & 0.936686340205 & 0.065406801766 & 100.000000000000 & 0.010000000000 & 1528.892979019463\\
5 & 0.962814033966 & 0.037894996979 & 1.000000000000 & 1.000000000000 & 26.388707737612\\
5 & 0.962814033966 & 0.037894996979 & 2.000000000000 & 0.500000000000 & 52.777415475225\\
5 & 0.962814033966 & 0.037894996979 & 10.000000000000 & 0.100000000000 & 263.887077376124\\
5 & 0.962814033966 & 0.037894996979 & 100.000000000000 & 0.010000000000 & 2638.870773761235\\
7 & 0.980707416582 & 0.019481114073 & 1.000000000000 & 1.000000000000 & 51.331766563503\\
7 & 0.980707416582 & 0.019481114073 & 2.000000000000 & 0.500000000000 & 102.663533127006\\
7 & 0.980707416582 & 0.019481114073 & 10.000000000000 & 0.100000000000 & 513.317665635029\\
7 & 0.980707416582 & 0.019481114073 & 100.000000000000 & 0.010000000000 & 5133.176656350294\\
11 & 0.992105208579 & 0.007926120285 & 1.000000000000 & 1.000000000000 & 126.165130478895\\
11 & 0.992105208579 & 0.007926120285 & 2.000000000000 & 0.500000000000 & 252.330260957791\\
11 & 0.992105208579 & 0.007926120285 & 10.000000000000 & 0.100000000000 & 1261.651304788953\\
11 & 0.992105208579 & 0.007926120285 & 100.000000000000 & 0.010000000000 & 12616.513047889535\\
13 & 0.994336053642 & 0.005680047328 & 1.000000000000 & 1.000000000000 & 176.054871067187\\
13 & 0.994336053642 & 0.005680047328 & 2.000000000000 & 0.500000000000 & 352.109742134375\\
13 & 0.994336053642 & 0.005680047328 & 10.000000000000 & 0.100000000000 & 1760.548710671875\\
13 & 0.994336053642 & 0.005680047328 & 100.000000000000 & 0.010000000000 & 17605.487106718749\\
\bottomrule
\end{tabular}
\end{table}

